\documentclass[12pt]{article}
\usepackage{amsmath, amssymb}
\usepackage{geometry}
\usepackage{booktabs}
\usepackage{parskip}

\geometry{margin=1in}

\title{Binomial Distribution\\{\large ORF 309 --- Based on Slides07}}
\author{}
\date{}

\begin{document}
\maketitle

%---------------------------------------------------------------
\section*{1. Start with a Bernoulli Random Variable}

A \textbf{Bernoulli RV} is any random variable that takes only the values $0$ or $1$.

\[
X_i = \begin{cases} 1 & \text{``success'' (e.g.\ heads), with probability } p \\ 0 & \text{``failure'' (e.g.\ tails), with probability } q = 1 - p \end{cases}
\]

A \textbf{Bernoulli Process} is an infinite sequence $X_1, X_2, X_3, \ldots$ of independent, identically distributed Bernoulli RVs, each with the same parameter $p$.

%---------------------------------------------------------------
\section*{2. Counting Successes: $S_n$}

Define the random variable
\[
S_n = X_1 + X_2 + \cdots + X_n,
\]
which counts the \textbf{total number of successes in $n$ trials}. $S_n$ takes values in $\{0, 1, 2, \ldots, n\}$.

%---------------------------------------------------------------
\section*{3. The Probability Formula: $P(S_n = k)$}

\textbf{Step 1.} Each specific sequence with exactly $k$ successes and $n - k$ failures has probability
\[
p^k \cdot (1-p)^{n-k}.
\]
For example, with $n = 8$, $k = 2$: the sequence $\mathtt{ssffffff}$ has probability $p^2(1-p)^6$.

\textbf{Step 2.} Every such sequence has the \emph{same} probability (just the positions of the $k$ successes rearranged).

\textbf{Step 3.} Count how many distinct arrangements exist:
\[
\binom{n}{k} = \frac{n!}{k!\,(n-k)!}
\qquad \text{(read: ``$n$ choose $k$'')}
\]
For $n=8$, $k=2$: $\binom{8}{2} = \frac{8 \times 7}{2 \times 1} = 28$ arrangements.

\textbf{Putting it together:}
\[
\boxed{P(S_n = k) = \binom{n}{k}\, p^k\,(1-p)^{n-k}}
\]

%---------------------------------------------------------------
\section*{4. Formal Definition}

A random variable $X$ with values $\{0, 1, 2, \ldots, n\}$ has the \textbf{Binomial Distribution} if
\[
P(X = k) = \binom{n}{k}\, p^k\,(1-p)^{n-k}, \qquad k = 0, 1, \ldots, n,
\]
where $p \in [0,1]$ and $q = 1-p$. We write $X \sim \mathrm{Binomial}(n, p)$.

\textbf{Sanity check} (probabilities sum to 1):
\[
\sum_{k=0}^{n} \binom{n}{k} p^k (1-p)^{n-k} = (p + (1-p))^n = 1^n = 1. \checkmark
\]

%---------------------------------------------------------------
\section*{5. Expected Value: $E[X] = np$}

\subsection*{Method 1 (long way)}

Plug the formula directly into $E[X] = \sum_{k=0}^{n} k \cdot P(X = k)$ and simplify algebraically. It works, but it is tedious.

\subsection*{Method 2 (indicator trick --- the one to use)}

Write $X$ as a sum of indicator RVs:
\[
X = X_1 + X_2 + \cdots + X_n.
\]
By \textbf{linearity of expectation} (no independence required):
\[
E[X] = E[X_1] + E[X_2] + \cdots + E[X_n].
\]
Each $E[X_i] = P(X_i = 1) = p$, so
\[
\boxed{E[X] = \underbrace{p + p + \cdots + p}_{n \text{ terms}} = np.}
\]

%---------------------------------------------------------------
\section*{6. Application to Q3a (HW3)}

\textbf{Setup.} You draw one coin and flip it \textbf{twice} independently. Let $X$ = number of heads.

Conditioned on the coin type, the flips are independent and each has the same $p$, so $X \mid \text{coin type} \sim \mathrm{Binomial}(2,\, p)$.

\begin{center}
\begin{tabular}{lccc}
\toprule
Condition & $n$ & $p$ & $E[X] = np$ \\
\midrule
Fair coin     & 2 & 0.50 & $2 \times 0.50 = \mathbf{1}$ \\
Canadian coin & 2 & 0.65 & $2 \times 0.65 = \mathbf{1.3}$ \\
\bottomrule
\end{tabular}
\end{center}

You do \emph{not} need to compute $P(X=0)$, $P(X=1)$, $P(X=2)$ individually. The indicator trick gives the answer immediately.

\end{document}
